\section*{अध्याय \ref{ch:b2:dispersion-wave-packets} --- परिक्षेपण और तरंग-पुंज}

\begin{solution}{exo:b2:dispersion-wave-packets:1}
(क) $k = \qty{0.063}{rad/m}$, $\omega = \sqrt{gk} = \qty{0.79}{rad/s}$, $T = \qty{8.0}{s}$;
$v_\varphi = \qty{12.5}{m/s}$, $v_g = \qty{6.3}{m/s}$। (ख) $\omega = \qty{0.45}{rad/s}$, $k = \omega^2
/g = \qty{0.021}{rad/m}$, $\lambda = \qty{310}{m}$; $v_\varphi = gT/2\pi = \qty{22}{m/s}$, $v_g =
\qty{11}{m/s}$; $5 \times 10^6/11 = \qty{4.6e5}{s} = 5.3$ दिन। (ग) $\lambda = \qty{6}{m}$, $v_g
= \qty{1.6}{m/s}$: उसे पाँच सप्ताह लगते, और छोटी तरंगें उससे बहुत पहले
\omterm{def:b2:viscous-flows:viscosity}{श्यानता} से अवमंदित और हवा से तोड़ दी जाती हैं।
\end{solution}

\begin{solution}{exo:b2:dispersion-wave-packets:2}
$\omega = \qty{1.26e8}{rad/s}$, $\omega_c = \qty{6.3e7}{rad/s}$: $k = \sqrt{\omega^2 - \omega_c^2}/c =
\qty{0.36}{rad/m}$, $v_\varphi = \qty{3.5e8}{m/s}$, $v_g = c^2k/\omega = \qty{2.6e8}{m/s}$,
और गुणनफल $c^2$। \qty{5}{MHz} पर $\omega < \omega_c$: $\kappa = \sqrt{\omega_c^2 - \omega^2}/c =
\qty{0.18}{m^{-1}}$, गहराई \qty{5.5}{m}।
\end{solution}

\begin{solution}{exo:b2:dispersion-wave-packets:3}
\qty{4}{Hz} पर विस्पंद, और प्रबलता का आवर्तकाल \qty{0.25}{s}। पानी: $f = 0.125$ और
\qty{0.119}{Hz}; आवरण $\cos(\delta\omega t - \delta kx)$, जिसमें $\delta k = \qty{2.9e-3}{rad/m}$
(अर्थात् \qty{1.1}{km} की दूरी पर विस्पंद) और चाल $\delta\omega/\delta k = \qty{6.4}{m/s}$ ---
अर्थात् \qty{105}{m} पर \omterm{prop:b2:dispersion-wave-packets:group}{समूह वेग}, $\tfrac12\sqrt{g\lambda/2\pi} = \qty{6.4}{m/s}$।
\end{solution}

\begin{solution}{exo:b2:dispersion-wave-packets:4}
$\ln(b/a) = 1.20$: $\Lambda = \qty{0.24}{\micro H/m}$, $\Gamma = \qty{106}{pF/m}$, $c = 1/
\sqrt{\Lambda\Gamma} = \qty{2.0e8}{m/s}$, $Z_c = \qty{48}{\ohm}$, \qty{5.0}{ns/m}; \qty{10}{ns}
$\leftrightarrow$ \qty{2.0}{m}। हवा: $\Gamma = \qty{46}{pF/m}$, $c = c_0$, $Z_c = \qty{72}{\ohm}$।
\end{solution}

\begin{solution}{exo:b2:dispersion-wave-packets:5}
(क) $\omega^2 = c^2k^2 + (EI/\mu)k^4$; $v_\varphi = \sqrt{c^2 + (EI/\mu)k^2}$; $v_g = (c^2k +
2(EI/\mu)k^3)/\omega$। (ख) $c^2 = \qty{1.12e5}{m^2/s^2}$, $k = \qty{66}{rad/m}$, $(EI/\mu)k^2
= \qty{7.0e3}{m^2/s^2}$: $v_\varphi/c = \sqrt{1.063} = 1.031$, अर्थात् $+3.1\%$ --- वही $53$ सेंट जो
\cref{ch:b2:waves-on-strings} में मिले थे। (ग) $v_g > v_\varphi$: असामान्य; ऊँचे
संनादी आगे निकल जाते हैं और छेड़ की शुरुआत किसी छोटी उतरती
चहक में बदल जाती है।
\end{solution}

\begin{solution}{exo:b2:dispersion-wave-packets:6}
(क) $-\mu\omega^2 + \iu\alpha\omega = -Tk^2$: $\underline k^2 = (\omega^2/c^2)(1 - \iu\alpha/\mu\omega)$।
(ख) $\underline k \approx (\omega/c)(1 - \iu\alpha/2\mu\omega)$: $k' = \omega/c$, $k'' = \alpha/2\mu c =
\alpha/2Z$; और $8.7k''$ dB/m। (ग) $c = \qty{100}{m/s}$, $Z = \qty{0.5}{kg/s}$, $\mu\omega = 3.1 \gg
0.02$: $k'' = \qty{0.02}{m^{-1}}$, $\delta = \qty{50}{m}$, \qty{0.17}{dB/m}। (घ) इस सीमा में
नहीं ($k' = \omega/c$); पर जब $\alpha \gtrsim \mu\omega$ हो तो वर्गमूल का
रैखिकीकरण नहीं होता और $k'$ $\omega$ पर अरैखिक रूप से निर्भर करने लगता है: अर्थात् परिक्षेपी (केबल की
विसरणशील सीमा)।
\end{solution}

\begin{solution}{exo:b2:dispersion-wave-packets:7}
(क) $v_\varphi = \sqrt{\gamma k/\rho}$, $v_g = \tfrac32v_\varphi$: अर्थात् समूह शिखरों से आगे
निकल जाता है। (ख) $v_\varphi^2 = g/k + \gamma k/\rho$ $k_m = \sqrt{\rho g/\gamma}$ पर न्यूनतम है:
$\lambda_m = 2\pi\sqrt{\gamma/\rho g} = \qty{1.7}{cm}$, $v_{\min} = (4g\gamma/\rho)^{1/4} =
\qty{23}{cm/s}$। (ग) अनुतरंग नाव के सापेक्ष अचल कोई आकृति है,
जो उन तरंगों से बनी है जिनका \omterm{def:b2:dispersion-wave-packets:relation}{कला वेग} उसकी चाल से मेल खाता है; और कोई भी
\qty{23}{cm/s} से धीमी नहीं। (घ) केशिका दशा में, $v_g > v_\varphi$: तब शिखर
छल्ले के अगले किनारे पर जन्मते हैं और उसके पिछले किनारे पर मरते हैं।
\end{solution}

\begin{solution}{exo:b2:dispersion-wave-packets:8}
(क) $k = \qty{0.126}{rad/m}$, $\omega'' = -\tfrac14\sqrt{g/k^3} = \qty{-18}{m^2/s}$: $t_{\text{फैलाव}}
= 2.5 \times 10^5/18 = \qty{1.4e4}{s}$, अर्थात् चार घंटे। (ख) \qty{10}{ps}: $t_{\text{फैलाव}} =
4 \times 10^{-6}/0.16 = \qty{25}{\micro s}$, और \qty{2e8}{m/s} पर \qty{5}{km}; \qty{1}{ns}: $0.04/0.16 =
\qty{0.25}{s}$, \qty{50000}{km}। छोटे स्पंद (ऊँची बिट-दरें) किलोमीटरों के भीतर
फैल जाते हैं: अतः लंबी कड़ियाँ प्रतिकारी तंतु या मध्यम दरें काम में लेती हैं।
\end{solution}

\begin{solution}{exo:b2:dispersion-wave-packets:9}
(क) $\partial_x^2v = R\Gamma\,\partial_tv$, $D = 1/R\Gamma = \qty{1.1e12}{m^2/s}$। (ख) $\underline k^2
= -\iu R\Gamma\omega$: $k' = k'' = \sqrt{R\Gamma\omega/2}$; \qty{1}{Hz}: $k = \qty{1.7e-6}{m^{-1}}$,
$v_\varphi = \qty{3.7e6}{m/s}$, $\delta = \qty{600}{km}$; \qty{10}{Hz}: $v_\varphi = \qty{1.2e7}{m/s}$,
$\delta = \qty{190}{km}$। (ग) प्रति बिंदु $R\Gamma L^2 = 9 \times 10^{-13} \times 9 \times 10^{12} = \qty{8}{s}$:
अर्थात् लगभग एक शब्द प्रति मिनट। (घ) $\Lambda = R\Gamma/G = \qty{9}{H/km}$ --- किसी
समुद्री केबल के लिए यह असंगत प्रेरकत्व है; भारण-कुंडलियाँ थल की
दूरभाष-रेखाओं पर चलीं, जिनके $G$ और $R\Gamma$ भिन्न थे।
\end{solution}

\begin{solution}{exo:b2:dispersion-wave-packets:10}
(क) $\partial_te = \dot y(\mu\ddot y + Ky) + Ty'\dot y' = \dot yTy'' + Ty'\dot y' = \partial_x
(Ty'\dot y) = -\partial_x\mathcal P$। (ख) $\langle e\rangle = \tfrac14A^2(\mu\omega^2 + Tk^2 + K) =
\tfrac12\mu\omega^2A^2$ ($\mu\omega^2 = Tk^2 + K$ लेकर); $\langle\mathcal P\rangle = \tfrac12Tk\omega
A^2$। (ग) अनुपात $Tk/\mu\omega = c^2k/\omega = v_g$। (घ) $y = A\eu^{-\kappa x}\cos\omega t$:
$\mathcal P \propto \sin2\omega t$, माध्य शून्य --- अर्थात् कोई ऊर्जा माध्यम को पार नहीं करती।
\end{solution}

\begin{solution}{exo:b2:dispersion-wave-packets:11}
(क) $\tanh kh \to 1$: $\omega^2 = gk$; और $kh \ll 1$: $\omega = \sqrt{gh}\,k$, चाल $\sqrt{gh}$,
अपरिक्षेपी। (ख) $kh = 0.13$: उथला; $c = \sqrt{9.81 \times 4000} = \qty{198}{m/s}$,
$T = \qty{17}{min}$; \qty{11}{h} में \qty{8000}{km}। (ग) $\tfrac12\rho gA^2 = \qty{1.2}{kJ/m^2}$,
और $200 \times 1000$ km$^2$ पर: अर्थात् \qty{2.5e14}{J}। ग्रीन: $A \propto h^{-1/4}$, $(4000/10)^{1/4}
= 4.5$: अतः \qty{2.2}{m}, खड़े होने और चढ़ाव से पहले। (घ) $h \approx \lambda/2 =
\qty{50}{m}$ पर; और शिखर का जो भाग उथले पानी में होता है वह धीमा पड़ जाता है, अतः
शिखर घूमकर गहराई-रेखाओं के समांतर हो जाता है --- अर्थात् अपवर्तन।
\end{solution}

\begin{solution}{exo:b2:dispersion-wave-packets:12}
(क) $x = \ell$ पर: $v = v_i + v_r$, $i = (v_i - v_r)/Z_c$, $v = R_Li$: अतः $r = (R_L -
Z_c)/(R_L + Z_c)$। (ख) $0$, $-1$, $+1$, $0.5$। (ग) रेखा का प्रवेश
$Z_c$ जैसा दिखता है, अतः दोलनदर्शी $t = 0$ पर \qty{0.5}{V} देखता है, और फिर
प्रतिध्वनि $0.5r$ V $2\ell/c = \qty{1}{\micro s}$ पर: अर्थात् कुछ नहीं, \qty{-0.5}{V}, \qty{+0.5}{V}, \qty{+0.25}{V}; और
मिलाया हुआ जनित्र प्रतिध्वनि सोख लेता है, अतः दूसरी कोई नहीं आती। (घ)
$\tfrac12 \times 2 \times 10^8 \times 1.8 \times 10^{-6} = \qty{180}{m}$।
\end{solution}

\begin{solution}{pb:b2:dispersion-wave-packets:1}
\textbf{1.} $\partial_xv = -\Lambda\partial_ti - Ri$, $\partial_xi = -\Gamma\partial_tv$; और
$\eu^{\iu(\omega t - kx)}$ लेकर: $\iu kv = (\iu\omega\Lambda + R)i$, $\iu ki = \iu\omega\Gamma v$, अतः
$k^2 = \omega\Gamma(\omega\Lambda - \iu R)$।

\textbf{2.} $\omega = R/\Lambda = \qty{4000}{rad/s}$, \qty{640}{Hz}: अतः कुछ हर्ट्ज़ पर $\Lambda\omega
\ll R$, अर्थात् विसरणशील।

\textbf{3.} $k' = k'' = \sqrt{R\Gamma\omega/2}$, जहाँ $R\Gamma = \qty{5e-13}{s/m^2}$: \qty{1}{Hz}:
$\qty{1.25e-6}{m^{-1}}$, $v_\varphi = \qty{5e6}{m/s}$, $\delta = \qty{800}{km}$; \qty{4}{Hz}:
$\qty{2.5e-6}{m^{-1}}$, $\qty{1e7}{m/s}$, \qty{400}{km}।

\textbf{4.} $8.7L/\delta$: \qty{1}{Hz} पर \qty{38}{dB}, \qty{4}{Hz} पर \qty{76}{dB}: अतः केवल
सबसे धीमे घटक बचते हैं; और कोई तीखा बिंदु किसी लंबे, नीचे कूबड़ की तरह आता है।

\textbf{5.} $R\Gamma L^2 = 5 \times 10^{-13} \times 1.2 \times 10^{13} = \qty{6}{s}$: अर्थात् हर छह
सेकंड में एक बिंदु, और एक शब्द प्रति मिनट (असली केबल ने संवेदनशील दर्पण-धारामापियों
के साथ कुछ अधिक निकाल लिए थे)।

\textbf{6.} $\Lambda = R/\omega = \qty{80}{mH/km}$। तब $k^2 = \Gamma R\omega(1 - \iu)$, $k = \sqrt{\Gamma R\omega}
\,2^{1/4}\eu^{-\iu\pi/8}$: $k' = \qty{3.9e-6}{m^{-1}}$, $k'' = \qty{1.6e-6}{m^{-1}}$: अतः $v_\varphi =
\qty{6.4e6}{m/s}$, $\delta = \qty{620}{km}$, और पूरी केबल पर \qty{76}{dB} के बदले
\qty{49}{dB}।

\textbf{7.} हर \qty{6}{s} में एक बिंदु लगभग \qty{0.2}{bit/s} है: अतः तंतु
$10^{10}$ से $10^{11}$ गुना तेज़ है।

\textbf{8.} $v_\varphi = gT/2\pi$, $v_g = gT/4\pi$।

\textbf{9.} \qty{18}{s}: $v_g = \qty{14.1}{m/s}$, $4 \times 10^6/14.1 = \qty{3.3}{d}$;
\qty{6}{s}: \qty{4.7}{m/s}, \num{9.9} दिन: अतः स्वेल छह दिन से अधिक चलती है।

\textbf{10.} $t = D/v_g = 4\pi D/gT$, अतः $1/T = gt/4\pi D$। \qty{86400}{s} में $\Delta(1/T) = 1/12 - 1/16
= \qty{0.0208}{Hz}$: अतः $g/4\pi D = \qty{2.4e-7}{s^{-2}}$, $D = \qty{3200}{km}$।

\textbf{11.} $\tfrac12\rho gA^2 = \qty{11}{kJ/m^2}$; और ऊर्जा $v_g = gT/4\pi =
\qty{9.4}{m/s}$ से चलती है: अतः प्रति मीटर शिखर \qty{103}{kW}।

\textbf{12.} $10^3 \times 103 \times 0.3 = \qty{31}{MW}$: अर्थात् दस पवन-चक्कियाँ।

\textbf{13.} $\lambda = 2\pi U^2/g = \qty{16}{m}$। हर तरंग-घटक की ऊर्जा
अपनी कला-चाल की आधी चाल से चलती है, अतः आकृति नाव के पीछे किसी फाँक में
सिमटी रहती है (अर्ध-कोण \qty{19.5}{\degree}, चाल चाहे कुछ भी हो)।

\textbf{14.} कोई शिखर समूह के पिछले सिरे पर जन्मता है, उससे आगे निकलता है
(शिखर पानी के सापेक्ष $2v_g$ से चलते हैं, और समूह के सापेक्ष $v_g$ से)
और अगले सिरे पर मर जाता है। समूह $120 \times 9.4 = \qty{1.1}{km}$
लंबा है, अर्थात् पाँच तरंगदैर्घ्य ($\lambda = gT^2/2\pi = \qty{225}{m}$); और प्लव से दस शिखर
गुज़रते हैं, जिनमें से हर एक लगभग दो मिनट जिया।

\textbf{15.} $kh = 0.13$: उथला। $c = \qty{198}{m/s}$, $T = \lambda/c = \qty{17}{min}$,
\qty{8.4}{h}।

\textbf{16.} अपरिक्षेपी: सारे घटक \qty{198}{m/s} पर, अतः वे एक ही स्पंद बनकर
साथ आते हैं (जबकि स्वेल, $\sim\qty{10}{m/s}$ पर, आवर्तकाल से छँटकर
एक सप्ताह लेती)।

\textbf{17.} $\tfrac12\rho gA^2 = \qty{1.8}{kJ/m^2}$, $500 \times 200$ km$^2$ पर: अर्थात् \qty{1.8e14}{J},
जो किसी मेगाटन का $4\%$ है।

\textbf{18.} $A = 0.6(4000/20)^{1/4} = \qty{2.3}{m}$; $c = \sqrt{gh} = \qty{14}{m/s}$, $\lambda =
cT = \qty{14}{km}$। शिखर, जो अपने आगे की गर्त से गहरे पानी में है,
तेज़ चलता है और उसे पकड़ लेता है: अतः अग्र-भाग खड़ा होकर ज्वार-भित्ति बन जाता है।

\textbf{19.} पंद्रह मिनट में \qty{0.6}{m} की चढ़ाई, और ढाल
$10^{-5}$: भाँपने को कुछ नहीं।

\textbf{20.} माध्य चाल, अर्थात् माध्य गहराई के लिए $\sqrt{gh}$; और चूँकि
$\sqrt h$ नतोदर है, उथले खंड तरंग को जितना धीमा करते हैं उतना गहरे
तेज़ नहीं करते, अतः पथ का औसत $\sqrt{g\bar h}$ से नीचे रहता है।

\textbf{21.} $\omega' = \tfrac12\sqrt{g/k}$, $\omega'' = -\tfrac14\sqrt{g/k^3}$; $k = \qty{0.063}{rad/m}$:
$\omega'' = \qty{-50}{m^2/s}$।

\textbf{22.} $t_{\text{फैलाव}} = 10^6/50 = \qty{2e4}{s}$, \qty{5.6}{h}; और पुंज
$v_gt = 6.25 \times 2 \times 10^4 = \qty{125}{km}$ चल चुका होता है।

\textbf{23.} लंबी रेल में $k$ की सँकरी पट्टी होती है, जिनके \omterm{prop:b2:dispersion-wave-packets:group}{समूह
वेग} लगभग मिलते हैं: अतः घटक अधिक देर साथ रहते हैं।
सुनामी का माध्यम अपरिक्षेपी है: बराबर चालें, कोई फैलाव नहीं।

\textbf{24.} $\omega'' = \hbar/m = \qty{1.2e-4}{m^2/s}$: $t_{\text{फैलाव}} = 10^{-18}/1.2 \times
10^{-4} = \qty{9e-15}{s}$; $v_g = \hbar k/m = \qty{7.3e6}{m/s}$: अर्थात् \qty{60}{nm}।

\textbf{25.} केबल: $k^2 = -\iu R\Gamma\omega$, $v_\varphi = \sqrt{2\omega/R\Gamma}$, कोई साफ़
$v_g$ नहीं, और वह क्षीण भी करती है तथा लीप भी देती है। स्वेल: $\omega^2 = gk$, $v_\varphi = 2v_g$,
फैलती है। सुनामी: $\omega = \sqrt{gh}\,k$, $v_\varphi = v_g$, दोनों में से कुछ नहीं। इलेक्ट्रॉन:
$\omega = \hbar k^2/2m$, $v_g = 2v_\varphi$, फैलता है।
\end{solution}
